\begin{figure}[!ht]
    \centering
    \includegraphics[width=\linewidth]{Images/Optimization.png}
\caption{\textbf{Optimization Process Enhancements}. The PIML problem involves satisfying multiple constraints in a given domain. For complex constraints, the problem can be simplified using sequential training or stacked training~\citep{howard2023stacked}. When dealing with complex domains (e.g., large, irregular), domain decomposition can be employed~\citep{shukla2021parallel}. (b) The loss function measures the cumulative error in the domain and, once discretized, comprises global weights~\citep{wang2021understanding}, local weights~\citep{anagnostopoulos2024residual}, a sampling method~\citep{wu2023comprehensive}, and a transformation function. (c) The optimizer updates the parameters using a line search method~\citep{urban2024unveiling}. The choice of the symmetric matrix defines the optimization method, such as gradient descent, Adam, L-BFGS, or Newton’s method. Several approaches have improved base performance, ensuring conflict-free updates~\citep{liu2024config}, or by using alternative methods like non-dominated genetic algorithms~\citep{akhter2024common}, or particle swarm optimization~\citep{davi2022pso}.}

    \label{Optim_Enh}
\end{figure}
